% Modernised reading — fonds Alexandre Grothendieck, shelfmark 75
% (pages 1-32, the whole folder).
%
% This is an INTERPRETATION, not a transcription. It re-reads the pages in
% today's notation and vocabulary, states the hypotheses the manuscript leaves
% implicit, and drops the critical apparatus so that the mathematics reads
% straight through. Everything the brackets and underlines carried moves into
% a footnote. It is held to being mathematically correct as it stands; where
% the manuscript is loose or mistaken, the true statement is what appears here
% and the footnote says what the page has. In any dispute of substance the
% transcription governs: whoever cites Grothendieck must cite batch-01.fr.tex
% (pp. 1-20) or batch-02.fr.tex (pp. 21-32).
%
% Pass: Opus 5.5 (claude-opus-5-5), 2026-10-10 — first pass, unchecked against the pages by a human.
% The folder was read whole, its two batches together; this is one document
% for the shelfmark. It was made on Opus 5.5 and has not been measured against
% the reference reading of folder 115 (made on Opus 5). The literature named
% in the footnotes is cited from memory and was not checked against the
% sources. Models followed: the readings of folders 81, 88 and 89 of the same
% group, cross-referenced where they meet this one.
%
% Scope. Page 1 is the cover (« C4 DEA — Géométrie combinatoire », « (24 p) »);
% its words head the spine section. Pages 2-27 are his hand, course
% preparation. Pages 28-29 are two copies of a typed letter of 30.01.78 to the
% Minister of Education, skipped by the transcription as carrying no
% mathematics; it is named here only as the transcription's header names it.
% Pages 30-32 are his own typed announcements of the two 1977-78 courses.
%
% Spine. The « yoga » of torsors, built from scratch and applied to polygons:
%   p. 2       homogeneous spaces, torsors, Aut of a torsor;
%   pp. 3-7    graphs, contours, orientations, u_omega, Aut(X) dihedral;
%   pp. 8-12   Z-sets, conjugacy classes in S_n, semi-direct products, Aff(Z),
%              the sections q_t and sigma_t, fibre products, pull-backs;
%   pp. 12-13  small symmetric groups as affine/projective groups over F_q;
%              vector/affine/projective structures as orbits of Isom;
%   pp. 14-15  regular realisations of polygons;
%   pp. 15-20  contracted product, (G,E)-structures, items (1)-(7): objects
%              isomorphic to E <-> Aut(E)-torsors, Env(P) = ^T C, functors on a
%              connected groupoid <-> G-objects, extension of the structure
%              group, the product of torsors under a commutative group;
%   p. 21      (8) « structures cubiques »: signed bases <-> sets with a
%              fixed-point-free involution;
%   pp. 22-23  a contour = (omega, torsor under omega ∧ Z_n); torsors under a
%              semi-direct product;
%   pp. 24-27  orientations as {±1}-torsors in three contexts; associativity
%              formulas; flags (announced);
%   pp. 30-32  the two course announcements.
%
% Notation fixed for the document. Torsors are RIGHT torsors, sets to be
% twisted are LEFT G-sets, and the contracted product is written
% P ∧_G E = (P x E)/G = ^P E throughout (the page writes × or ∧, both).
% Z_n = Z/nZ; D_n = Z_n ⋊ {±1} (order 2n); semi-direct products written
% Z ⋊ Γ (he writes Z.Γ, the normal factor first). GL, PGL for his Gl, GP, with
% GP(m,k) = PGL(m+1,k) (his GP(m,k) acts on P^m(k)). omega_X is the
% two-element set of orientations of X, a {±1}-torsor; u_omega the elementary
% rotation of p. 7. S_n, A_n gothic.
%
% Substantive departures from the pages, each footnoted where it occurs:
% — p. 5 a) and p. 7: the two conditions 1°) (s adjacent to gs) and 2°) (a
%   adjacent to ga) are called equivalent; they are for n = 3 and n >= 5 and
%   for the infinite contour, NOT for the square (n = 4), where the two
%   reflections through edge midpoints satisfy 1° only and the two diagonal
%   reflections 2° only. The reading takes the conjunction. Same finding as
%   folder 89, p. 20.
% — p. 5, Corollary: « Z/2Z . Z/H » puts the normal factor second, against
%   his own convention of p. 8; written Z/H ⋊ Z/2Z.
% — p. 6: « groupe cyclique d'ordre n >= 2 » leaves out the infinite contour
%   (G^0 = Z); supplied.
% — p. 8: « u transitif <=> tous les n_i sauf un sont nuls » needs « and that
%   one equals 1 »; corrected. Table completed (S_5: 7 classes; A_4: 4; A_5:
%   5) — the completion is ours.
% — p. 10: i: Z -> Aut(Z,T) and q_{tg} = int(i(g)) q_t need Z commutative
%   (true in every application: Z_n, V); stated.
% — p. 13: A_4 ≃ Aff(1,F_4) « ? » is true; said. P(V) = V/k^* corrected to
%   (V - {0})/k^*. The projective dimension (n in the sentence, n-1 in the
%   formula) regularised.
% — p. 17: « homothétie par zeta^i » is the complex homothety, i.e. the real
%   rotation; and Env(P) carries the regular realisation of circumradius 1,
%   not the unimodular one.
% — p. 18: « G -> Aut(X)^° » for a LEFT action is G -> Aut(X); the ° (opposite)
%   is read as a vestige of the cancelled « droite ». Same on p. 20.
% — p. 19: « E G-ens. (à gauche), E x_G G' » mixes sides; written with T a
%   right torsor, T' = T ∧_G G', F a left G'-set.
% — p. 23: the Exer holds as stated for Z commutative; for Z arbitrary the
%   group of T_0 is T ∧_G Z (conjugation), isomorphic to omega ∧_Γ Z only up
%   to inner automorphisms.
% — pp. 24-25: b) needs k connected (else det lands in mu_2(k), not {±1});
%   « P/Gl+ », « P/SO », « P/A_n » equal P ∧_G {±1} only when s is onto
%   (n >= 1, resp. n >= 2 for sets); the reading keeps P ∧_G {±1} as the
%   definition, which the associativity formulas need. The a)-b) diagram
%   with sgn det makes sense for k = R only.
% — p. 26: the concatenated basis is written over I^- only; it is over all of
%   I, I^- first (or any total order inducing the given one on I^-, margin NB).
%
% Modern names given and footnoted as ours: holomorph, hyperoctahedral group
% (Weyl group of type B_n), twisted form, Picard groupoid, heap (Prüfer, Baer),
% graded determinant (Knudsen-Mumford), exceptional isomorphisms, PSL/PGL.
%
% Announced and never established in the folder: the topological realisation
% (p. 3, cancelled), the « Dénombrement » of contours (p. 4), the description
% of non-connected contours (p. 6), S_5 and A_4 in the table (p. 8), the
% (G,E)-structure <-> torsor statement promised on p. 15 (made only as (2) on
% p. 17), item (7) (cancelled), the « Compatibilité entre les isom. can.
% a) b) c) » (p. 27) and the « Cas d'extensions » (p. 27, « Définir »).

\documentclass[11pt,a4paper]{article}
\input{../preamble/grothendieck}

\folder{75}
\pages{1}{32}
\dating{1977-1978}
\watermark{Édition de démonstration\\interprétation personnelle de l'œuvre}
\foldertitle{Cours C4, DEA : notes manuscrites (s.d.), lettre (1978), tapuscrit (1977-1978) — lecture modernisée du dossier entier}

\begin{document}

\section*{Résumé}

\begin{resume}
Ces feuillets préparent deux cours que Grothendieck donne à Montpellier en
1977--1978 : un cours « C4 » sur la géométrie de l'icosaèdre et un cours
de DEA sur les polygones réguliers « au-dessus d'un anneau de base
quelconque ». Le dossier se ferme sur les deux annonces dactylographiées de ces
cours ; tout ce qui précède est l'outillage qu'il juge nécessaire pour les
faire, écrit pour lui-même, vite.

Le problème de départ est d'une grande simplicité. Un polygone à $n$ côtés,
dessiné sur une feuille, a des sommets, des côtés, et deux sens de parcours. Si
on oublie le dessin et qu'on ne garde que la question « quel sommet touche quel
côté ? », il reste un objet combinatoire --- un \emph{contour} --- dont on veut
comprendre les symétries, les orientations, et les façons de le réaliser comme
une vraie figure, dans le plan réel, mais aussi sur un corps fini ou un anneau.

L'idée qui organise tout est celle de \emph{torseur}. Un cadran d'horloge sans
chiffres est un bon exemple : on peut y faire tourner l'aiguille d'un, deux,
trois crans, mais aucune position n'est « midi ». Le groupe des rotations agit
sans qu'aucun point soit privilégié. Un torseur, c'est cela : un ensemble sur
lequel un groupe agit comme sur lui-même, mais sans origine. Les sommets d'un
polygone orienté forment un torseur sous le groupe cyclique ; les deux
orientations d'un polygone forment un torseur sous le groupe $\{\pm 1\}$ ; les
façons de numéroter les sommets forment un torseur sous le groupe diédral.
Grothendieck en tire un principe qu'il appelle lui-même un « yoga » : se donner
un objet ayant la forme d'un objet modèle $E$, c'est exactement se donner un
torseur sous le groupe des symétries de $E$. Tout polygone à $n$ côtés est le
polygone standard « tordu » par un tel torseur, et toute construction naturelle
faite sur le polygone standard --- par exemple sa réalisation régulière dans le
plan complexe --- se transporte d'elle-même à tous les autres. On obtient ainsi
pour chaque polygone combinatoire un plan euclidien canonique où il est
régulier, sans avoir jamais choisi ni origine ni sens.

Le même principe donne un calcul des orientations. Orienter un espace
vectoriel réel, numéroter à parité près un ensemble fini, orienter un module
quadratique : chaque fois l'ensemble des deux orientations est un torseur sous
$\{\pm 1\}$, et ces torseurs se multiplient. Il écrit la règle qui donne
l'orientation d'une somme directe à partir de celles des facteurs --- avec un
terme de correction qui ne dépend que de l'ordre des facteurs de dimension
impaire --- et la même règle pour un ensemble découpé en morceaux.

En chemin, les petits groupes symétriques apparaissent comme des groupes de
géométrie sur les corps finis : les permutations de quatre objets sont les
transformations projectives d'une droite sur le corps à trois éléments, et
celles de cinq objets sur le corps à cinq éléments. Ce sont ces
« isomorphismes exceptionnels » que l'annonce du cours sur l'icosaèdre met au
centre.

Les deux annonces disent aussi quelque chose de sa façon d'enseigner. Il
s'étonne que l'icosaèdre, étudié depuis plus de deux mille ans, n'ait jamais
été regardé de façon systématique comme objet combinatoire ; il propose de
partir « de l'intuition combinatoire la plus naïve », sur des modèles en carton
aux arêtes coloriées, et appelle ses cours une « investigation commune », dont
les fruits dépendent du travail personnel de l'étudiant. Ce qu'on voit dans les
feuillets va dans le même sens : il ne fixe jamais un choix quand il peut
garder tous les choix ensemble et faire apparaître ce qui n'en dépend pas.

\keywords{torsor, homogeneous space, semidirect product, dihedral group, cycle graph, contracted product, twisted form, groupoid, orientation torsor, hyperoctahedral group, exceptional isomorphisms, regular polygon, icosahedron}
\end{resume}

\pagerange{1}{32}
\section*{Le fil du dossier, et les conventions}

La chemise (page 1) porte au crayon « C4 DEA --- Géométrie combinatoire », et
au-dessus « (24 p) »\footnote{Le compte ne correspond pas aux trente-deux pages
du dossier ; il pourrait ne compter que les feuillets manuscrits des pages 2 à
27, à quelques unités près, mais rien sur la chemise ne le dit. « Géométrie
combinatoire » est une lecture à vérifier.}. Ici les deux cours, « C4 » et « DEA », sont ceux que décrivent les
pages 30 à 32. Les feuillets n'ont pas de pagination de sa main ; à partir de
la page 16, il numérote les articles d'une série de (1) à (8), le (8) ouvrant
la page 21. Plus haut, les pages 8 à 12 ont leur propre série 1) à 4).

\subsection*{Les stations}

\begin{enumerate}
\item[1.] \emph{Rappels sur groupes, espaces homogènes, torseurs} (page 2).
\item[2.] \emph{Graphes et contours} (pages 3 à 7) : graphes, réalisation
rectiligne, contours, orientations, la rotation élémentaire $u_\omega$, le
groupe des automorphismes d'un contour.
\item[3.] \emph{Ensembles à permutation, produits semi-directs} (pages 8 à
12) : classes de conjugaison de $\mathfrak{S}_n$, groupe affine, sections
$q_t$ et réflexions $\sigma_t$ d'un polygone, produits fibrés, images
inverses d'extensions.
\item[4.] \emph{Petits groupes et géométrie sur les corps finis} (pages 12 et
13), puis structures vectorielle, affine, projective comme orbites.
\item[5.] \emph{Polygones réguliers} (pages 14 et 15), en exercice.
\item[6.] \emph{Le yoga des torseurs} (pages 15 à 19) : produit contracté,
$(G,E)$-structures, articles (1) à (4), l'enveloppe d'un polygone, les
foncteurs sur un groupoïde connexe.
\item[7.] \emph{Extension du groupe structural et produit des torseurs}
(pages 19 et 20), articles (5) à (7).
\item[8.] \emph{Structures cubiques} (page 21), article (8).
\item[9.] \emph{Un contour comme torseur tordu} (pages 22 et 23).
\item[10.] \emph{Orientations et formules d'associativité} (pages 24 à 27).
\item[11.] \emph{La lettre et les deux annonces de cours} (pages 28 à 32).
\end{enumerate}

Le dossier ne recommence pas sur lui-même. Il dépose deux fois la même
matière sous deux angles : les contours des pages 3 à 7 sont décrits à la main,
par sommets et arêtes, puis aux pages 18 et 22 par des torseurs, et ce sont
les mêmes objets. Il en va de même des orientations des pages 5 à 7, définies
sur le graphe, et de celles des pages 24 à 27, définies par un torseur et un
caractère.

\subsection*{Conventions}

Un \emph{torseur} sous un groupe $G$ est un ensemble non vide $P$ muni d'une
action \emph{à droite} de $G$, libre et transitive ; les ensembles qu'on tord
sont des $G$-ensembles \emph{à gauche}, et le produit contracté est noté
\[
  P \wedge_G E = (P \times E)/\bigl((pg, x) \sim (p, gx)\bigr) = {}^{P}E
\]
partout, là où la page écrit tantôt $\times_G$, tantôt $\wedge_G$\footnote{Il
écrit aussi une fois $P_E$ et une fois $F^{T}$ pour ${}^{P}E$, ${}^{T}F$. Les
côtés des actions ne sont pas toujours tenus sur la page ; on les signale là
où cela change un énoncé.}. On note $\mathbf{Z}_n = \mathbf{Z}/n\mathbf{Z}$,
$\mathbf{D}_n = \mathbf{Z}_n \rtimes \{\pm 1\}$ le groupe diédral d'ordre
$2n$, et $Z \rtimes \Gamma$ le produit semi-direct de $\Gamma$ par le
sous-groupe distingué $Z$, qu'il écrit $Z.\Gamma$. On écrit $\mathrm{GL}$ et
$\mathrm{PGL}$ pour ses $\mathrm{Gl}$ et $\mathrm{GP}$, en notant que son
$\mathrm{GP}(m, k)$ est le groupe projectif de $\mathbf{P}^{m}(k)$,
c'est-à-dire $\mathrm{PGL}(m+1, k)$. Pour un objet $X$, $\omega_X$ est
l'ensemble à deux éléments de ses orientations, torseur sous $\{\pm 1\}$, et
$\omega \wedge \omega'$ le produit contracté de deux tels torseurs.

\subsection*{Ce que le dossier annonce sans l'établir}

La réalisation topologique d'un graphe (page 3, esquissée puis biffée) ; le
« dénombrement » des contours (page 4) ; la description des contours non
connexes (page 6) ; les lignes de $\mathfrak{S}_5$ et de $\mathfrak{A}_4$ du
tableau de la page 8 ; l'équivalence entre $(G,E)$-structures et torseurs
promise page 15 (elle est énoncée, sans plus, comme article (2) de la page 17) ;
l'article (7) sur les orientations, biffé page 20 et repris page 24 ; la
« compatibilité entre les isomorphismes canoniques a), b), c) » et le « cas
des extensions » de la page 27, que la page se donne à faire.

\pagerange{2}{2}
\section*{1. Groupes, espaces homogènes, torseurs (page 2)}

Un groupe $G$ opère sur un ensemble $E$ ; l'opération est \emph{transitive}
si deux points quelconques sont échangés par un élément de $G$, et $E$ est un
\emph{espace homogène} s'il est de plus non vide. Les couples $(E, x)$ d'un
espace homogène et d'un point correspondent aux sous-groupes $H$ de $G$ (le
stabilisateur de $x$, et inversement $E = G/H$, $x = H$), et c'est une
équivalence de catégories : il existe un morphisme $(E, x) \to (E', x')$
--- forcément unique, un morphisme d'espaces homogènes étant déterminé par
l'image d'un point --- si et seulement si $H \subset H'$. Sans point marqué :
$E$ et $E'$ sont isomorphes si et seulement si $H$ et $H'$ sont conjugués, et
il existe un morphisme $E \to E'$ si et seulement si $H$ est conjugué à un
sous-groupe de $H'$. Les classes d'isomorphisme d'espaces homogènes
correspondent donc aux classes de conjugaison de sous-groupes\footnote{Cette
dernière phrase est écrite en diagonale dans la marge ; « $\Leftrightarrow$ »
et « conjugaison » y sont des lectures.}.

Un \emph{torseur} est un espace homogène dont les stabilisateurs sont réduits à
l'élément neutre ; de façon équivalente, un ensemble non vide sur lequel, pour
tous $x, y$, il existe un unique $g$ avec $y = x g$ ; de façon équivalente
encore, un $G$-ensemble isomorphe à $G$ opérant sur lui-même par translations.
Les torseurs sont tous isomorphes, non canoniquement, et ce sont les « plus
gros » espaces homogènes : tout espace homogène en est un quotient.

\textbf{Automorphismes d'un torseur.} Le groupe $\mathrm{Aut}_G(P)$ des
automorphismes d'un torseur $P$ est isomorphe à $G$, non canoniquement. Avec
nos conventions ($P$ torseur à droite), le choix de $x \in P$ définit
$u_x : G \xrightarrow{\sim} \mathrm{Aut}_G(P)$ par $u_x(g)(x h) = x g h$, et si
l'on remplace $x$ par $y = x s$, on trouve
\[
  u_{xs} = u_x \circ \mathrm{int}(s), \qquad \mathrm{int}(s)(g) = s g s^{-1}.
\]
\footnote{La page fait opérer $G$ à gauche ($y = s x$) et écrit, encadré,
$u_{sx} = u_x \circ \mathrm{int}(s^{-1})$, après une vérification où
figure un point d'interrogation ; c'est la même formule du côté gauche. Avec
les conventions du document : $u_{xs}(g)$ envoie $x s h$ sur $x s g h$, et
$u_x(s g s^{-1})$ envoie $x s h$ sur $x (s g s^{-1}) s h = x s g h$.}
Le groupe des automorphismes d'un torseur est donc canoniquement $G$ à
automorphisme intérieur près --- canoniquement $G$ si $G$ est commutatif.
C'est le fait qui reviendra aux pages 10 et 23.

\pagerange{3}{7}
\section*{2. Graphes et contours (pages 3 à 7)}

\pagerange{3}{4}
\subsection*{Graphes et réalisations (pages 3 et 4)}

Un \emph{graphe} est la donnée d'un ensemble $S$ de sommets, d'un ensemble $A$
d'arêtes et d'une relation d'incidence $R \subset S \times A$, avec un seul
axiome : chaque arête a exactement deux sommets, $\operatorname{card} R^{-1}(a) =
2$\footnote{Il écrit $R\{a\}$ pour l'ensemble des sommets de $a$, et plus
loin $R\{s\}$ pour l'ensemble des arêtes en $s$. L'axiome exclut les boucles
mais permet plusieurs arêtes entre deux mêmes sommets, comme sur le dessin de
la page (deux arêtes doubles, une arête et un sommet isolés).}. Le graphe est
\emph{strict} si une arête est déterminée par ses sommets ; un graphe strict
est donc un ensemble $S$ muni d'une partie $A$ de $\mathfrak{P}_2(S)$, l'ensemble
des parties à deux éléments --- un graphe simple, dans le vocabulaire
d'aujourd'hui. Un \emph{arc}, ou arête orientée, est une arête munie du choix
d'une de ses extrémités comme origine.

\emph{Réalisation topologique.} À une partie $a$ à deux éléments on associe le
segment
\[
  \operatorname{seg}(a) = \Bigl\{ (t_x)_{x \in a} \in [0,1]^{a} \Bigm| \textstyle\sum_x t_x = 1 \Bigr\},
\]
qui contient $a$ comme ensemble de ses deux extrémités ; la réalisation du
graphe est le recollement
\[
  \operatorname{réal}(S, A, R) = \Bigl( S \sqcup \coprod_{a \in A}
  \operatorname{seg} R^{-1}(a) \Bigr) \Big/ \sim
\]
où chaque extrémité d'un segment est identifiée au sommet correspondant de
$S$\footnote{La page esquisse aussi une version par les arcs,
$(S \sqcup (\vec{A} \times [0,1]))/\!\sim$, biffée, et décrit le recollement
comme celui de $X = \coprod \operatorname{seg}$ le long de
$\partial X \to S$ ; « recollement » y est une lecture et un mot est
illisible. On munit la réalisation de la topologie quotient.}.

\emph{Réalisation rectiligne.} Pour un graphe fini, une réalisation rectiligne
dans un espace vectoriel réel $E$ de dimension finie est un plongement
$\operatorname{réal}(S, A, R) \hookrightarrow E$ envoyant chaque arête sur un
segment.

\textbf{Proposition} (page 4). Un graphe fini admet une réalisation
rectiligne si et seulement s'il est strict, et il en admet alors une dans
$\mathbf{R}^3$. Plus précisément, si $S \hookrightarrow E \simeq \mathbf{R}^3$
est une injection telle que trois points de $S$ ne soient jamais alignés et
que, pour $x, y, z, t$ distincts dans $S$, les segments $[x,y]$ et $[z,t]$ soient
disjoints, elle se prolonge de façon unique en une réalisation rectiligne.

La nécessité est claire : deux arêtes de mêmes extrémités auraient la même
image. Pour l'existence, il suffit de placer les sommets en position générale,
quatre jamais coplanaires --- par exemple sur la courbe $t \mapsto (t, t^2,
t^3)$ ---, ce qui assure les deux conditions ; l'unicité tient à ce qu'une
arête n'a qu'une image possible, le segment joignant ses extrémités\footnote{Les
précisions « ($\simeq \mathbf{R}^3$) », « distincts », « i.e. trois points
jamais alignés » et la condition sur $x, y, z, t$ sont ajoutées au crayon ;
l'argument par la courbe des moments est de nous. Dans $\mathbf{R}^2$
l'énoncé serait faux : un graphe fini strict n'a de réalisation rectiligne
plane que s'il est planaire, ce qui exclut déjà le graphe complet à cinq
sommets.}.

Un graphe est \emph{connexe} si sa réalisation l'est, c'est-à-dire s'il n'est
pas somme --- réunion disjointe, $S \sqcup S'$, $A \sqcup A'$,
$R \sqcup R'$ --- de deux graphes non vides\footnote{La ligne « graphe connexe
$=$ réal. géom. connexe $=$ quasi-\ldots » se termine sur une lecture
incertaine ; on la complète par la ligne suivante de la page, qui parle des
sommes de deux graphes non vides.}. Tout graphe est, de façon unique, réunion
disjointe de graphes connexes.

\pagerange{4}{5}
\subsection*{Contours et leurs automorphismes (pages 4 et 5)}

Un \emph{contour} (combinatoire) --- il dit aussi graphe polygonal ou
multipolygonal --- est un graphe dont chaque sommet appartient à exactement
deux arêtes. Les composantes connexes d'un contour sont des contours, et
\textbf{un contour connexe non vide est déterminé à isomorphisme près par le
cardinal $c$ de $S$}, qui est aussi celui de $A$ : $2 \leq c \leq \aleph_0$.
Les contours connexes correspondent aux sous-groupes $H \neq \mathbf{Z}$ de
$\mathbf{Z}$, avec $c = \operatorname{card} \mathbf{Z}/H$ : pour $H =
n\mathbf{Z}$, $n \geq 2$, le polygone à $n$ côtés, le \emph{digone} pour
$n = 2$, et pour $H = 0$ la droite infinie, dont les sommets sont les entiers.
Le contour est strict si et seulement si $c \geq 3$\footnote{Ce dernier point
est au crayon. La page annonce aussi un « dénombrement » des contours,
lecture incertaine, qu'elle ne fait pas.}.

Soit $X = (S, A, R)$ un contour connexe strict ($c \geq 3$). Disons que deux
sommets sont \emph{liés} s'ils sont les extrémités d'une même arête, et deux
arêtes liées si elles ont un sommet commun. Pour un automorphisme $g$ de $X$,
considérons les conditions
\begin{enumerate}
\item[1°)] pour tout $s \in S$, $s$ est lié à $g s$ ;
\item[2°)] pour tout $a \in A$, $a$ est lié à $g a$.
\end{enumerate}
Un automorphisme qui satisfait 1°) et 2°) sera appelé \emph{rotation
élémentaire} de $X$\footnote{La page déclare 1°) et 2°) « conditions
équivalentes » et nomme $g$ « autom. d'orientation ». L'équivalence est vraie
pour $c = 3$, pour $c \geq 5$ et pour le contour infini, mais fausse pour le
carré : les deux réflexions par les médiatrices de deux côtés opposés
envoient chaque sommet sur un voisin mais fixent deux arêtes, et les deux
réflexions par les diagonales envoient chaque arête sur une arête voisine
mais fixent deux sommets. Les rotations d'un cran sont exactement les
automorphismes qui satisfont les deux conditions ensemble ; c'est ce qu'on
retient. Le dossier 89 (page 20) rencontre la même difficulté pour le carré,
sur la seule condition 1°).}.

\textbf{Proposition} (page 5). Soit $\omega(X)$ l'ensemble des rotations
élémentaires de $X$.
\begin{enumerate}
\item[b)] $\omega(X)$ a deux éléments, inverses l'un de l'autre :
$\omega(X) = \{g, g^{-1}\}$.
\item[c)] Le groupe $\mathrm{Aut}(X)$ opère par conjugaison sur $\omega(X)$, et
le morphisme $\mathrm{Aut}(X) \to \mathfrak{S}_{\omega(X)} \simeq
\mathbf{Z}/2\mathbf{Z}$ est surjectif ; son noyau $\mathrm{Aut}^{+}(X)$ est un
sous-groupe d'indice $2$, cyclique, engendré par l'un quelconque des $g \in
\omega(X)$.
\end{enumerate}
\textbf{Corollaire.} $\mathrm{Aut}(X) \simeq \mathbf{Z}/H \rtimes
\mathbf{Z}/2\mathbf{Z}$, où $H \subset \mathbf{Z}$ est le sous-groupe des $m$
tels que $g^m = 1$ : le groupe diédral $\mathbf{D}_n$ si $c = n$ est fini, le
groupe diédral infini si $c = \aleph_0$\footnote{La page écrit
« $\operatorname{Aut}(X) \simeq \mathbf{Z}/2\mathbf{Z} \cdot \mathbf{Z}/H$
(produits ½ directs) », le facteur distingué en second, contrairement à sa
propre convention de la page 8 ($Z.\Gamma$, $Z$ distingué). Le « tt » de
« engendré par tt $g$ » est une lecture.}.

\pagerange{6}{7}
\subsection*{Orientations, et la rotation élémentaire $u_\omega$ (pages 6 et 7)}

Une \emph{orientation} d'un contour $(S, A, R)$ est le choix, pour toute arête
$a$, de l'une de ses extrémités $\sigma(a)$ comme origine, de sorte que
\[
  \text{(O)} \qquad \text{pour tout } s \in S, \text{ il existe une unique
  arête } a \ni s \text{ avec } \sigma(a) = s,
\]
autrement dit que $\sigma : A \to S$ soit bijective. Si $\omega$ est une
orientation, l'orientation opposée $-\omega$ choisit partout l'autre
extrémité.

\textbf{Proposition} (page 7). Un contour connexe a exactement deux
orientations, $\omega$ et $-\omega$. Pour une orientation $\omega$, soit
$u_\omega(s)$ l'autre extrémité de l'arête $\sigma^{-1}(s)$ issue de $s$, et
$u_\omega(a)$ l'arête issue de l'extrémité de $a$. Alors $u_\omega$ est un
automorphisme du graphe, $u_{-\omega} = u_\omega^{-1}$, et :
\begin{enumerate}
\item[(i)] l'ordre de $u_\omega$ est $n = \operatorname{card} S =
\operatorname{card} A$ (infini si $n$ l'est) ; donc $u_\omega \neq u_{-\omega}$
si et seulement si $n \geq 3$ ;
\item[(ii)] pour $n \geq 3$, $\omega \mapsto u_\omega$ est une bijection de
l'ensemble $\omega_X$ des orientations sur l'ensemble $\omega(X)$ des rotations
élémentaires de la page 5 ;
\item[(iii)] le groupe cyclique engendré par $u_\omega$ est le groupe $G^0 =
\mathrm{Aut}^{+}(X)$ des automorphismes qui respectent l'orientation, et $S$
et $A$ sont des torseurs sous $G^0$.
\end{enumerate}
On identifiera désormais l'ensemble $\omega_X$ des deux orientations et
l'ensemble $\omega(X)$ des deux rotations élémentaires\footnote{La page 5
définit $\omega(X)$ comme une partie de $\mathrm{Aut}(X)$, la page 7 nomme
$\mathrm{Or}(S,A,R)$ l'ensemble des orientations et établit la bijection ; à
la page 22, $\omega(C)$ est l'ensemble des orientations. Sur (ii), la page
renvoie aux mêmes « conditions équivalentes » 1°, 2° qu'à la page 5 ; voir la
note correspondante pour le carré.}. Un torseur sous $\mathbf{Z}_n$ est un
ensemble à $n$ éléments muni d'un ordre circulaire, et un contour orienté
connexe à $n$ sommets n'est rien d'autre.

\textbf{Descriptions d'un contour connexe} (page 6). Se donner un contour
orienté connexe, c'est se donner un entier $n \geq 2$, ou $n = \infty$, et un
torseur sous $\mathbf{Z}_n$ (sous $\mathbf{Z}$ si $n = \infty$) : l'ensemble
$S$ des sommets, l'arête issue de $s$ joignant $s$ à $s + 1$. Sans orientation,
se donner un contour connexe de $n \geq 3$ sommets revient à se donner
\begin{enumerate}
\item[1°)] un groupe cyclique $G^0$, d'ordre $n$ (ou monogène infini) ;
\item[2°)] une paire $\{g, g^{-1}\}$ de générateurs inverses l'un de
l'autre ;
\item[3°)] un torseur $S$ sous $G^0$,
\end{enumerate}
les arêtes étant les paires $\{s, s g\}$. Pour $n = 2$ la paire $\{g, g^{-1}\}$
se réduit à un point, et l'on remplace 2°) par
\begin{enumerate}
\item[2 bis)] un ensemble $\Omega$ à deux éléments $\{a, b\}$ et une application
$i : \Omega \to G^0$ avec $i(a)\,i(b) = 1$, $i(a)$ engendrant $G^0$,
\end{enumerate}
qui convient pour tout $n$ : $\Omega$ est l'ensemble des orientations, et $i$
l'application $\omega \mapsto u_\omega$\footnote{La paire $\{g,
g^{-1}\}$ est la « bi-circulation » $\{u, u^{-1}\}$ du dossier 81 (pages 41
et 42), et la troisième description d'un polygone du dossier 89 (pages 19 et
20). La page dit « groupe cyclique
$G^0$ (d'ordre $n \geq 2$) », ce qui laisse de côté le contour infini ; on
l'ajoute. L'interprétation de $\Omega$ et de $i$ est de nous ; elle est celle
que la page 7 justifie. La phrase qui commence l'article à la page 6
(« \ldots{} $A$ --- choix --- NB $\sigma : S \to A$ est compatible avec
l'action de $G^0$ ») est à moitié illisible ; une phrase de la fin de la
page 7, « se donner un contour », s'arrête au bas du feuillet.}.

Enfin les contours orientés non nécessairement connexes sont les ensembles sur
lesquels $\mathbf{Z}$ opère \emph{sans point fixe} --- c'est la
« rectification » de la page 8 : une orbite à un élément serait un polygone à
un côté, une boucle, que l'axiome des graphes interdit. La description
annoncée des contours non orientés non connexes n'est pas écrite.

\pagerange{8}{12}
\section*{3. Ensembles à permutation et produits semi-directs (pages 8 à 12)}

\pagerange{8}{8}
\subsection*{Ensembles à permutation (page 8)}

Un ensemble muni d'une permutation $u$ est un $\mathbf{Z}$-ensemble. Ses
orbites sont des espaces homogènes $\mathbf{Z}/i\mathbf{Z}$, $i \in
\mathbf{N}$ ($i = 0$ pour les orbites infinies), et le $\mathbf{Z}$-ensemble
est classé à isomorphisme près par la famille de cardinaux $(n_i)_{i \in
\mathbf{N}}$, $n_i$ étant le nombre d'orbites isomorphes à
$\mathbf{Z}/i\mathbf{Z}$. L'ensemble est fini si et seulement si $n_0 = 0$ et
si les $n_i$ sont finis et presque tous nuls ; $u$ est transitif si et
seulement si $\sum_i n_i = 1$, c'est-à-dire si tous les $n_i$ sont nuls sauf
un, qui vaut $1$\footnote{La page écrit « tous les $n_i$ sauf un sont nuls »,
sans « qui vaut 1 » ; deux mots illisibles suivent, dont le premier est lu
« exemplaire ». Elle note la famille $(a_i)$, puis $(n_i)$.}. Pour $E$ fini
non vide, $u$ transitif revient à un ordre circulaire sur $E$.

\textbf{Corollaire.} Les classes de conjugaison de $\mathfrak{S}_n$
correspondent aux partitions de $n$, familles $(c_i)_{1 \leq i \leq n}$
d'entiers avec $\sum i\,c_i = n$ ($c_i$ cycles de longueur $i$). Ainsi :
\[
  \begin{array}{lll|lll}
    \mathfrak{S}_2 & 2 & 1,\ (12) &
      \mathfrak{A}_2 & 1 & 1 \\
    \mathfrak{S}_3 & 3 & 1,\ (12),\ (123) &
      \mathfrak{A}_3 & 3 & 1,\ (123),\ (132) \\
    \mathfrak{S}_4 & 5 & 1,\ (12),\ (12)(34),\ (123),\ (1234) &
      \mathfrak{A}_4 & 4 & 1,\ (12)(34),\ (123),\ (132) \\
    \mathfrak{S}_5 & 7 & &
      \mathfrak{A}_5 & 5 &
  \end{array}
\]
\footnote{La page s'arrête à « $\mathfrak{S}_5$ : ? » et laisse la ligne de
$\mathfrak{A}_4$ en points de suspension ; une première ligne de
$\mathfrak{A}_4$, « $\sigma_{12}\sigma_{34}$, $\sigma_{1234}$ », est raturée
et entourée --- à juste titre, $(1234)$ étant impaire. Les nombres de classes
de $\mathfrak{S}_5$, $\mathfrak{A}_4$ et $\mathfrak{A}_5$ et les représentants
de $\mathfrak{A}_4$ sont de nous : dans $\mathfrak{A}_4$ les $3$-cycles se
répartissent en deux classes de quatre éléments. Il note $\sigma_{12}$,
$\sigma_{123}$, $\sigma_{321}$ les cycles écrits ici $(12)$, $(123)$,
$(132)$.}

\textbf{Rectification.} Les contours orientés sont les ensembles à
permutation $(E, u)$ où $u$ n'a pas de point fixe ; voir la page 6.

\pagerange{8}{9}
\subsection*{Produits semi-directs et groupe affine (pages 8 et 9)}

Une suite exacte $1 \to Z \xrightarrow{i} G \xrightarrow{p} \Gamma \to 1$
munie d'une section $q$ de $p$ qui est un morphisme de groupes revient à la
donnée d'un sous-groupe $\Gamma' \subset G$ que $p$ envoie isomorphiquement sur
$\Gamma$, ou encore à celle d'une action $\Gamma \to \mathrm{Aut}(Z)$ ; on
reconstitue $G = Z \rtimes \Gamma$ comme l'ensemble $Z \times \Gamma$ muni de
la loi
\[
  (z, \gamma)(z', \gamma') = \bigl(z\, \gamma(z'),\ \gamma\gamma'\bigr).
\]
Les sous-groupes de $G$ contenant $Z$ correspondent aux sous-groupes $\Gamma'$
de $\Gamma$, et s'écrivent $Z \rtimes \Gamma'$\footnote{Cette remarque est en
marge. Une seconde note marginale, en petits caractères obliques, renvoie à la
théorie des extensions et à $H^2(\Gamma, Z)$ ; elle est presque entièrement
illisible et on ne la reconstitue pas. La page écrit d'abord la loi sous une
autre forme, biffée.}.

\textbf{Le groupe affine.} Pour un groupe $Z$, soit $Z_d$ l'ensemble $Z$ sur
lequel $Z$ opère par translations. Le groupe affine $\mathrm{Aff}(Z)$ est le
groupe des couples $(\gamma, \rho)$, $\gamma \in \mathrm{Aut}(Z)$, $\rho$
bijection de $Z$, tels que $\rho(g x) = \gamma(g)\rho(x)$ pour tous $g, x$ ---
les automorphismes de $Z_d$ comme ensemble à groupe d'opérateurs. On a une
suite exacte scindée
\[
  1 \to Z \xrightarrow{i} \mathrm{Aff}(Z) \underset{q}{\overset{p}{\rightleftarrows}}
  \mathrm{Aut}(Z) \to 1, \qquad q(\gamma) = (\gamma, \gamma),\quad i(g) =
  (\mathrm{id}, \tau_g),
\]
où $\tau_g$ est la translation par $g$ ; on a
$\mathrm{int}(q(\gamma))|_Z = \gamma$, et donc $\mathrm{Aff}(Z) \simeq Z
\rtimes \mathrm{Aut}(Z)$\footnote{C'est ce qu'on appelle l'\emph{holomorphe}
de $Z$ ; le mot n'est pas sur la page.}. Pour $\Gamma \subset
\mathrm{Aut}(Z)$, $\mathrm{Aff}^{\Gamma}(Z) = Z \rtimes \Gamma$ en est un
sous-groupe.

Un morphisme d'ensembles à opérateurs à groupe variable $(E, G) \to (E', G')$
est un couple $(\rho, \gamma)$ avec $\rho(g x) = \gamma(g)\rho(x)$ ; si $E$ est
un torseur, $\rho$ détermine $\gamma$.

Pour $Z$ le groupe additif d'un espace vectoriel $V$ sur un corps $R$, et
$\Gamma = \mathrm{GL}_R(V)$, on trouve $\mathrm{Aff}_R(V) \simeq V \rtimes
\mathrm{GL}_R(V)$, avec $(z, \gamma)(z', \gamma') = (z + \gamma(z'),
\gamma\gamma')$ ; si $V$ porte une forme quadratique non dégénérée $Q$, le
groupe des déplacements --- les transformations affines dont la partie linéaire
conserve $Q$ --- est $\mathrm{Dépl}(V, Q) \simeq V \rtimes \mathrm{O}(V,
Q)$\footnote{« Non dégénérée » est une lecture. La page dit « transf. affines
qui conservent la distance », ce qui n'a de sens que pour $R = \mathbf{R}$ et
$Q$ définie positive ; dans ce cas une application qui conserve la distance
est automatiquement affine.}.

\pagerange{10}{11}
\subsection*{Les sections $q_t$ et les réflexions d'un polygone (pages 10 et 11)}

Soit $Z$ un groupe \emph{commutatif} et $T$ un torseur sous $Z$. Soit
$\mathrm{Aut}(Z, T)$ le groupe des couples $(\gamma, \rho)$, $\gamma \in
\mathrm{Aut}(Z)$, $\rho$ bijection de $T$ avec $\rho(t g) = \rho(t)\gamma(g)$.
On a une suite exacte
\[
  1 \to Z \xrightarrow{i} \mathrm{Aut}(Z, T) \xrightarrow{p} \mathrm{Aut}(Z)
  \to 1, \qquad i(g) = (\mathrm{id}, t \mapsto t g),
\]
et chaque $t \in T$ définit une section $q_t$ : $q_t(\gamma)$ est l'unique
élément au-dessus de $\gamma$ qui fixe $t$. Si $t' = t g$, alors
\[
  q_{t'} = \mathrm{int}(i(g)) \circ q_t,
\]
car $i(g)\,q_t(\gamma)\,i(g)^{-1}$ est au-dessus de $\gamma$ et envoie $t'$ sur
$t g = t'$\footnote{La page n'écrit pas que $Z$ est commutatif. Il le faut
pour que les translations $t \mapsto t g$ soient des automorphismes au-dessus
de l'identité, c'est-à-dire pour que $i$ soit défini sans choix ; pour $Z$
quelconque le noyau de $p$ est $\mathrm{Aut}_Z(T)$, isomorphe à $Z$
seulement à automorphisme intérieur près (page 2). Toutes les applications
du dossier ont $Z$ commutatif. « À droite », pour le torseur, est une
lecture.}. Même chose pour le sous-groupe $\mathrm{Aut}^{\Gamma}(Z, T)$ des
couples avec $\gamma \in \Gamma$.

\emph{Exemple.} $Z = \mathbf{Z}_n$ : $\mathrm{End}(Z) = \mathbf{Z}_n$,
$\mathrm{Aut}(Z) = \mathbf{Z}_n^{*} \supset \{\pm 1\}$, et
$\mathrm{Aff}(\mathbf{Z}_n) = \mathbf{Z}_n \rtimes \mathbf{Z}_n^{*}$ contient
$\mathrm{Aff}^{\{\pm 1\}}(\mathbf{Z}_n) = \mathbf{Z}_n \rtimes \{\pm 1\} =
\mathbf{D}_n$, le groupe diédral d'ordre $2n$ --- pour $n \geq 3$, car pour
$n = 2$ on a $+1 = -1$ dans $\mathbf{Z}_2^{*}$\footnote{« Groupe diédral
d'indice $n$ » : « indice » est une lecture ; $n$ est l'indice du polygone,
l'ordre du groupe est $2n$.}.

Soit maintenant $T$ un torseur sous un groupe cyclique $Z$ d'ordre $n \geq 3$
muni d'une paire $\{u, u^{-1}\}$ de générateurs inverses --- c'est, par la
page 6, un polygone à $n$ côtés --- et la suite exacte
\[
  1 \to Z \to \mathrm{Aut}^{\pm 1}(Z, T) \xrightarrow{p} \{\pm 1\} \to 1 .
\]
Une section est déterminée par $\sigma = q(-1)$, un élément avec $p(\sigma) =
-1$ et $\sigma^2 = 1$ : une réflexion du polygone. Pour $t \in T$, soit
$\sigma_t = q_t(-1)$, l'unique réflexion qui fixe le sommet $t$. En écrivant
$g$ pour $i(g)$ et en utilisant $\sigma_t g \sigma_t^{-1} = g^{-1}$ :
\[
  \boxed{\ \sigma_{tg} = g\,\sigma_t\,g^{-1} = g^2 \sigma_t = \sigma_t g^{-2}\ }
\]
\textbf{Corollaire.} Si $n$ est impair, $g \mapsto g^2$ est bijectif et $t
\mapsto \sigma_t$ est une bijection des sommets sur les $n$ réflexions : à
chaque sommet correspond l'unique automorphisme non orienté qui le fixe. Si
$n$ est pair, les fibres de $t \mapsto \sigma_t$ ont deux éléments, un sommet
et le sommet \emph{opposé} ($\sigma_t = \sigma_{t'}$ si et seulement si $t' =
t g$ avec $g^2 = 1$), et l'image est formée des $n/2$ réflexions qui passent
par des sommets ; les $n/2$ autres passent par des milieux de
côtés\footnote{La dernière phrase est de nous. La page s'arrête à « ses fibres
sont toutes de cardinal 2… », puis « Sommet opposé, symétrie canonique » ; une
note marginale, très incomplète, semble décrire $\sigma_t = \sigma_{t'}$ par
l'appartenance de l'écart entre $t$ et $t'$ à $\{1, s\}$, $s$ l'élément
d'ordre $2$.}.

\pagerange{11}{12}
\subsection*{Produits fibrés et images inverses d'extensions (pages 11 et 12)}

Le produit fibré de $p : E \to S$ et $q : F \to S$ est $P = E \times_S F =
\{(x, y) \mid p x = q y\}$, avec ses projections $p' : P \to F$ et $q' : P \to
E$ ; c'est la solution d'un problème universel dans la catégorie des ensembles,
et, si les données sont des groupes et des morphismes, dans celle des groupes.
On a alors $\mathrm{Ker}\, p' \simeq \mathrm{Ker}\, p$ et $\mathrm{Ker}\, q'
\simeq \mathrm{Ker}\, q$, et si $p$ est surjectif (bijectif), $p'$ l'est
aussi. C'est le langage des carrés cartésiens.

Une extension $1 \to Z \to G \to \Gamma \to 1$ et un morphisme $\Gamma' \to
\Gamma$ donnent par image inverse l'extension $1 \to Z \to G \times_\Gamma
\Gamma' \to \Gamma' \to 1$, scindée si la première l'est.
\textbf{Remarque.} Toute extension scindée $Z \rtimes \Gamma$ provient par
image inverse, le long de l'action $\Gamma \to \mathrm{Aut}(Z)$, de
l'extension scindée $Z \rtimes \mathrm{Aut}(Z) = \mathrm{Aff}(Z)$.

\pagerange{12}{13}
\section*{4. Petits groupes et géométrie sur les corps finis (pages 12 et 13)}

Un ensemble à deux éléments est la même chose qu'un torseur sous
$\mathbf{Z}_2$, ou qu'une droite affine sur $\mathbf{F}_2$ --- tout ensemble à
deux éléments en porte une seule structure. Les petits groupes symétriques et
alternés se lisent ainsi comme groupes de géométrie sur les corps finis, en
notant $\mathrm{Aff}(m, k)$ le groupe affine de $k^m$ et $\mathrm{PGL}(m+1,
k)$ le groupe projectif de $\mathbf{P}^m(k)$ :
\begin{align*}
  \mathfrak{S}_2 &\simeq \mathrm{Aff}(1, \mathbf{F}_2), & \mathfrak{A}_2 &= 1, \\
  \mathfrak{S}_3 &\simeq \mathrm{Aff}(1, \mathbf{F}_3) = \mathbf{F}_3 \rtimes \mathbf{F}_3^{*}
    \simeq \mathrm{GL}(2, \mathbf{F}_2) = \mathrm{PGL}(2, \mathbf{F}_2), &
    \mathfrak{A}_3 &\simeq \mathbf{F}_3, \\
  \mathfrak{S}_4 &\simeq \mathrm{Aff}(2, \mathbf{F}_2) \simeq \mathrm{PGL}(2, \mathbf{F}_3), &
    \mathfrak{A}_4 &\simeq \mathrm{Aff}(1, \mathbf{F}_4), \\
  \mathfrak{S}_5 &\simeq \mathrm{PGL}(2, \mathbf{F}_5), &
    \mathfrak{A}_5 &\simeq \mathrm{PGL}(2, \mathbf{F}_4) = \mathrm{SL}(2, \mathbf{F}_4).
\end{align*}
Presque chaque fois le groupe de géométrie opère fidèlement sur un ensemble
ayant le bon nombre de points --- les $3$ points de $\mathbf{F}_3$ ou de
$\mathbf{P}^1(\mathbf{F}_2)$, les $4$ points de $\mathbf{F}_2^2$, de
$\mathbf{P}^1(\mathbf{F}_3)$ ou de $\mathbf{F}_4$, les $5$ points de
$\mathbf{P}^1(\mathbf{F}_4)$ --- et les ordres coïncident ; seul
$\mathrm{PGL}(2, \mathbf{F}_5)$, d'ordre $120$, opère sur les six points de
$\mathbf{P}^1(\mathbf{F}_5)$ et non sur cinq objets visibles\footnote{La page écrit $\mathrm{GP}(1,
\cdot)$ pour $\mathrm{PGL}(2, \cdot)$, et accompagne $\mathfrak{A}_4 \simeq
\mathrm{Aff}(1, \mathbf{F}_4)$ d'un point d'interrogation. L'isomorphisme est
vrai : $\mathrm{Aff}(1, \mathbf{F}_4)$, d'ordre $12$, opère fidèlement sur les
quatre points de $\mathbf{F}_4$, donc s'identifie à un sous-groupe d'indice
$2$ de $\mathfrak{S}_4$, qui ne peut être que $\mathfrak{A}_4$. Les
isomorphismes $\mathfrak{S}_5 \simeq \mathrm{PGL}(2, \mathbf{F}_5)$ et
$\mathfrak{A}_5 \simeq \mathrm{PGL}(2, \mathbf{F}_4)$ sont les
« isomorphismes exceptionnels » que reprend l'annonce du cours de la
page 31. La page ne dit pas sur quels cinq objets $\mathrm{PGL}(2,
\mathbf{F}_5)$ opère, et on ne le reconstitue pas.}.

\emph{Groupes classiques.} Sur un corps $k$, les suites
\[
  1 \to k^{*} \to \mathrm{GL}(V) \to \mathrm{PGL}(V) \to 1, \qquad
  1 \to k^{*} \to \mathrm{GL}(n, k) \to \mathrm{PGL}(n, k) \to 1
\]
ne sont pas scindées en général, et $\mathrm{PGL}(V)$ opère sur l'espace
projectif $\mathbf{P}(V) = (V \smallsetminus \{0\})/k^{*}$\footnote{La marge
écrit $\mathbf{P}(V) \simeq V/k^{*}$ ; il faut ôter $0$.}.

\emph{Structures comme orbites.} Une structure de $k$-espace vectoriel de
dimension $n$ sur un ensemble $V$ est un élément de
$\mathrm{Isom}_{\mathrm{Ens}}(k^n, V)/\mathrm{GL}(n, k)$ ; une structure
d'espace affine de dimension $n$ sur $E$, un élément de
$\mathrm{Isom}_{\mathrm{Ens}}(k^n, E)/\mathrm{Aff}(n, k)$ ; une structure
d'espace projectif de dimension $m$ sur $P$, un élément de
$\mathrm{Isom}_{\mathrm{Ens}}(\mathbf{P}^m(k), P)/\mathrm{PGL}(m+1,
k)$\footnote{La page parle d'espace projectif « de dim. $n$ » et écrit la
formule avec $\mathbf{P}^{n-1}(k)$ et $\mathrm{GP}(n-1, k)$, qui décrit la
dimension $n - 1$ ; on a fixé une seule lettre. Elle écrit aussi $V$ pour
l'ensemble qui porte la structure affine.}. Les inclusions $\mathrm{GL}(n, k)
\subset \mathrm{Aff}(n, k) \subset \mathrm{PGL}(n+1, k)$ traduisent les
passages : espace vectoriel de dimension $n$ $\mapsto$ (oubli de l'origine)
espace affine de dimension $n$ $\mapsto$ espace projectif de dimension $n$,
projectivisé de l'espace vectoriel enveloppe de dimension $n + 1$, dont
l'espace affine est le complémentaire d'un hyperplan.

Entre crochets, la page esquisse une définition des espaces affines sans
groupe donné d'avance : une structure sur un ensemble $E$ non vide est la
donnée d'un sous-groupe $\Gamma \subset \mathfrak{S}_E$ qui opère
simplement transitivement --- $E$ devient un torseur sous $\Gamma$ ---, et une
structure d'espace affine sur $k$ est la donnée d'un tel $\Gamma$ commutatif
et d'un morphisme d'anneaux unitaires $k \to \mathrm{End}_{\mathbf{Z}}(\Gamma)$,
qui fait de $\Gamma$ l'espace vectoriel des translations\footnote{Le passage
est rapide et en partie biffé ; la condition sur $\Gamma$ après « tel que »
n'est pas écrite, et « simplement transitivement » est ce que demande la
suite --- « ens. à groupe d'op. fidèle » est lisible, et ne suffit pas. Un
torseur sous un groupe non précisé est ce qu'on appelle un \emph{tas}
(\emph{heap}, Prüfer 1924, Baer 1929) ; la page n'emploie pas ce mot et ne cite
personne.}.

\pagerange{14}{15}
\section*{5. Polygones réguliers (pages 14 et 15)}

\emph{Exercice.} Soit $E$ un plan affine euclidien. Une réalisation
rectiligne d'un polygone combinatoire $(S, A, R)$ dans $E$ est \emph{régulière}
si tous les côtés ont même longueur et tous les angles intérieurs même mesure,
\emph{unimodulaire} si cette longueur est $1$\footnote{« Et tous les angles
itou » est ajouté dans une bulle, et lu avec doute. Les deux conditions sont
nécessaires dès $n = 4$ : le losange a ses côtés égaux, le rectangle ses
angles égaux. Des angles intérieurs égaux valent $(n-2)\pi/n < \pi$, ce qui
force la convexité ; une réalisation étant un plongement, les polygones
étoilés sont exclus.}.
\begin{enumerate}
\item[a)] Existence, par les racines de $z^n = 1$ dans $\mathbf{C}$.
\item[b)] Si $(V \supset P \supset S)$ et $(V' \supset P' \supset S')$ sont deux
$n$-gones, $n \geq 3$, réalisés régulièrement et unimodulairement, tout
isomorphisme combinatoire $(S, A, R) \to (S', A', R')$ est induit par une
\emph{unique} isométrie de $V$ sur $V'$\footnote{La page écrit d'abord « il
existe une similitude (et une seule) qui transforme $P$ en $P'$ », puis
précise que l'unicité porte sur la similitude --- isométrie dans le cas
unimodulaire, ajouté au-dessus --- qui induit un isomorphisme combinatoire
donné. Il y en a $2n$ en tout, une par élément de $\mathbf{D}_n$.}.
\item[c)] L'enveloppe affine d'un polygone combinatoire réalisé (celle de
l'ensemble de ses sommets) et son centre de gravité ; une orientation du
polygone définit une orientation de son plan enveloppe.
\end{enumerate}
La page 15 reprend : un polygone régulier peut se définir par les côtés et
les angles égaux, ou par la transitivité de son groupe de similitudes ---
sur les drapeaux, couples d'un sommet et d'un côté qui le contient\footnote{La
page écrit « transitivité groupe des similitudes… » sans dire sur quoi. Sur
les sommets seuls, le rectangle conviendrait ; sur les côtés seuls, le
losange.} ; existence, et unicité à similitude près, tout isomorphisme
combinatoire étant induit par une unique similitude.

\pagerange{15}{19}
\section*{6. Le yoga des torseurs (pages 15 à 19)}

\pagerange{15}{16}
\subsection*{Produit contracté et transport de structure (pages 15 et 16)}

Soit $P$ un $G$-torseur (à droite) et $E$ un $G$-ensemble (à gauche), et
$F = P \wedge_G E = {}^{P}E$. Chaque $p \in P$ donne une bijection $E \to F$,
$x \mapsto [p, x]$, d'où une application
\[
  P \to \mathrm{Bij}(E, F),
\]
compatible aux actions à droite de $G$ ($G$ opérant sur $\mathrm{Bij}(E, F)$
par $u \cdot g = u \circ g$). Si $G$ opère fidèlement sur $E$, il opère
librement sur $\mathrm{Bij}(E, F)$, et dans ce cas, et dans ce cas seulement,
$P \to \mathrm{Bij}(E, F)$ est injective ; son image est une orbite de $G$
dans $\mathrm{Bij}(E, F)$. Appelons \emph{$(G, E)$-structure} sur un ensemble
$F$ une telle orbite. Toute structure sur $E$ invariante par $G$ --- de groupe,
d'anneau, de module sur un anneau fixé, d'espace topologique --- se transporte
alors à $F$ par l'une quelconque des bijections de l'orbite, et le résultat ne
dépend pas de celle qu'on choisit ; de même pour une structure mixte sur une
famille d'ensembles $(E_i)$, comme un anneau et un module, qui se transporte
aux ${}^{P}E_i$\footnote{La fin de la page 15, en petits caractères au bas du
feuillet, n'est lue qu'en partie ; elle annonce que, $G$ opérant fidèlement,
une $(G,E)$-structure « équivaut à » un torseur sous $G$, ce que l'article (2)
énonce. La notion de $(G,E)$-structure est celle des $G$-structures
d'Ehresmann, à ceci près qu'ici rien n'est topologique ; le nom n'est pas sur
la page.}.

\pagerange{16}{17}
\subsection*{Les articles (1) à (3) (pages 16 et 17)}

\textbf{(1) Le yoga.} Soit $\mathcal{C}$ une espèce de structures, $E$ une
structure d'espèce $\mathcal{C}$ et $G = \mathrm{Aut}_{\mathcal{C}}(E)$. Les
structures isomorphes à $E$ correspondent aux torseurs sous $G$ :
\[
  F \longmapsto \mathrm{Isom}_{\mathcal{C}}(E, F), \qquad
  P \longmapsto P \wedge_G E .
\]
$\mathrm{Isom}(E, F)$ est un $G$-torseur à droite par $u \mapsto u \circ g$ ;
l'évaluation $\mathrm{Isom}(E, F) \times E \to F$, $(u, x) \mapsto u(x)$,
vérifie $(u g)(x) = u(g x)$ et induit un isomorphisme $\mathrm{Isom}(E, F)
\wedge_G E \xrightarrow{\sim} F$ ; dans l'autre sens, $u \mapsto i_u$, $i_u(x) =
[u, x]$, est un isomorphisme $P \xrightarrow{\sim} \mathrm{Isom}(E, P \wedge_G
E)$ de torseurs, car $i_{ug} = i_u \circ g$. On obtient une équivalence entre
le groupoïde $\mathcal{C}_E$ des structures isomorphes à $E$ et le groupoïde
$\mathrm{Tors}(G)$ des $G$-torseurs\footnote{Il écrit « les (str) données $E$
$\longleftrightarrow$ torseurs sous $G$ », « str » étant une lecture. En
termes d'aujourd'hui, les objets de $\mathcal{C}_E$ sont les \emph{formes} de
$E$ --- sur un ensemble de base, elles sont toutes triviales, et ce qui
compte est la fonctorialité ; au-dessus d'une base non ponctuelle, c'est
l'énoncé qui classe les formes tordues par $H^1(\,\cdot\,, G)$. Le mot
« torseur » est celui de ses exposés de géométrie algébrique ; on disait
avant « espace principal homogène ».}.

\textbf{(2)} Si $E$ est un $G$-ensemble fidèle, fixé une fois pour toutes,
se donner un $G$-torseur revient à se donner un ensemble $F$ muni d'une
$(G, E)$-structure. C'est un cas particulier de (1), en prenant pour
$\mathcal{C}$ l'espèce des $(G, E)$-structures, dont le groupe
d'automorphismes en $E$ est $G$.

\textbf{(3)} Inversement, pour une espèce $\mathcal{C}$ et $G =
\mathrm{Aut}_{\mathcal{C}}(E)$, se donner sur un ensemble $F$ une structure
d'espèce $\mathcal{C}$ \emph{isomorphe} à $E$ équivaut à se donner une
$(G, E)$-structure. Ainsi (1) devient un cas particulier de (2), et les deux
énoncés sont le même.

\pagerange{17}{19}
\subsection*{(4) L'enveloppe d'un polygone combinatoire (pages 17 à 19)}

Soit $P_n$ le polygone combinatoire standard de sommets $\mathbf{Z}_n$,
$n \geq 3$, et sa réalisation régulière $i \mapsto \zeta^i$, $\zeta =
\exp(2 i\pi/n)$, dans le plan euclidien $\mathbf{C}$\footnote{Il faut $n \geq
3$ pour que ce soit une réalisation rectiligne, comme le dit la page. Le mot lu
« Arcs » après « régulière » est incertain.}. Le groupe $\mathbf{D}_n =
\mathrm{Aut}(P_n)$ opère sur $\mathbf{C}$ par isométries : la rotation $\tau_i$
par multiplication par $\zeta^i$, la réflexion $\sigma : i \mapsto -i$ par
$z \mapsto \bar{z}$\footnote{La page dit « homothétie par $\zeta^i$ » : c'est
l'homothétie complexe, c'est-à-dire la rotation réelle d'angle
$2i\pi/n$.}, et l'injection $S(P_n) = \mathbf{Z}_n \hookrightarrow \mathbf{C}$ est
équivariante.

Par (1), se donner un polygone combinatoire $P$ à $n$ sommets revient à se
donner le $\mathbf{D}_n$-torseur $T = \mathrm{Isom}(P_n, P)$, et $P =
{}^{T}P_n$. On pose
\[
  \mathrm{Env}(P) = {}^{T}\mathbf{C} = T \wedge_{\mathbf{D}_n} \mathbf{C},
\]
et l'injection équivariante donne
\[
  S(P) = {}^{T}S(P_n) \hookrightarrow {}^{T}\mathbf{C} = \mathrm{Env}(P).
\]
Comme $\mathbf{D}_n$ opère par isométries, $\mathrm{Env}(P)$ est un plan
euclidien, et $P$ y est réalisé régulièrement, sans qu'on ait choisi ni sommet
de départ ni sens : chaque polygone combinatoire a un plan enveloppe canonique,
où il est régulier, défini à isométrie unique près\footnote{La réalisation
obtenue a ses sommets sur le cercle unité ; elle n'est pas unimodulaire au
sens de la page 14, ses côtés valant $2\sin(\pi/n)$. « Liée aux
$\mathbf{D}_n$-opérations » est une lecture.}.

\textbf{Le principe général} (page 18). Soit $G$ un groupe. Le groupoïde
$\mathrm{Tors}(G)$ est équivalent au groupoïde à un seul objet de groupe
d'automorphismes $G$ --- l'objet étant le torseur trivial $G_d$, dont les
automorphismes sont les translations à gauche\footnote{La page écrit
$\underline{\mathrm{Hom}}(\mathbf{e}, G) \overset{\approx}{\to}
\underline{\mathrm{Tors}}(G)$ ; on lit ainsi cette ligne, sans pouvoir assurer
le sens qu'il donne à $\mathbf{e}$.}. Pour toute catégorie $\mathcal{E}$,
\[
  \underline{\mathrm{Hom}}\bigl(\underline{\mathrm{Tors}}(G), \mathcal{E}\bigr)
  \xrightarrow{\ \approx\ } G\text{-}(\mathcal{E}), \qquad \Phi \longmapsto
  \Phi(G_d),
\]
où $G\text{-}(\mathcal{E})$ est la catégorie des objets $X$ de $\mathcal{E}$
munis d'une action à gauche de $G$, c'est-à-dire d'un morphisme $G \to
\mathrm{Aut}(X)$ ; un quasi-inverse est $X \mapsto (P \mapsto {}^{P}X)$, quand
les produits contractés existent dans $\mathcal{E}$\footnote{La page écrit
« $G \to \mathrm{Aut}(X)^{\circ}$ », après avoir corrigé « droite » en
« gauche ». Une action à gauche est un morphisme $G \to \mathrm{Aut}(X)$ ; le
groupe opposé conviendrait à une action à droite, et on lit le « $\circ$ »
comme une trace de la première rédaction. Le foncteur $\Phi(G_d)$ reçoit bien
une action à gauche, par fonctorialité, des automorphismes de $G_d$, qui sont
les translations à gauche.}.

Pour un objet $E$ d'une catégorie $\mathcal{C}$, de groupe $G =
\mathrm{Aut}_{\mathcal{C}}(E)$, le groupoïde $\mathcal{C}_E$ des objets
isomorphes à $E$ est équivalent à $\mathrm{Tors}(G)$, par $F \mapsto
\mathrm{Isom}(E, F)$. D'où le triangle d'équivalences

\begin{tikzcd}[column sep=large]
  \underline{\mathrm{Hom}}(\underline{\mathrm{Tors}}(G), \mathcal{E})
  \arrow[r, "\approx"] \arrow[d, "\approx"'] &
  G\text{-}(\mathcal{E}) \\
  \underline{\mathrm{Hom}}(\mathcal{C}_E, \mathcal{E})
  \arrow[ur, "{\Phi \mapsto \Phi(E)}"'] &
\end{tikzcd}

\noindent et la conclusion : \emph{la catégorie des foncteurs de
$\mathcal{C}_E$ dans $\mathcal{E}$ est équivalente, de façon canonique à
isomorphisme unique près, à celle des $G$-objets de $\mathcal{E}$}, par $\Phi
\mapsto \Phi(E)$, de quasi-inverse $X \mapsto \bigl(F \mapsto
{}^{\mathrm{Isom}(E, F)}X\bigr)$. La marge, sous la mention « DEA », le dit
pour tout groupoïde connexe $\mathcal{C}_0$ et tout objet $E_0$, avec $G =
\mathrm{Aut}(E_0)$\footnote{La flèche oblique de la page, dessinée à double
pointe et surchargée, va de $G\text{-}(\mathcal{E})$ vers
$\underline{\mathrm{Hom}}(\mathcal{C}_E, \mathcal{E})$, avec à côté « $\Phi
\leftrightarrow \Phi(E)$ » ; on l'oriente dans le sens de $\Phi \mapsto
\Phi(E)$. La marge écrit $E$ pour $E_0$.}.

Appliqué aux polygones (page 19) : se donner, pour tout polygone
combinatoire $P$ à $n$ sommets, un espace vectoriel réel $E(P)$, « de façon
compatible avec le transport de structure » --- un foncteur sur le groupoïde
des $n$-gones ---, c'est se donner le seul espace $E = E(P_n)$ avec l'action
de $\mathbf{D}_n$ ; on reconstitue $E(P) = {}^{T}E$, $T = \mathrm{Isom}(P_n,
P)$. L'enveloppe $\mathrm{Env}$ est l'exemple de $\mathbf{C}$.

\pagerange{19}{20}
\section*{7. Extension du groupe structural, produit des torseurs (pages 19 et 20)}

\textbf{(5) Extension du groupe d'opérateurs.} Soit $G \to G'$ un morphisme
de groupes et $T$ un $G$-torseur. Alors $T' = T \wedge_G G'$ est un
$G'$-torseur, et pour tout $G'$-ensemble à gauche $F$ --- donc $G$-ensemble par
restriction ---
\[
  T' \wedge_{G'} F = (T \wedge_G G') \wedge_{G'} F \simeq T \wedge_G (G'
  \wedge_{G'} F) \simeq T \wedge_G F, \qquad\text{soit}\quad {}^{T'}F \simeq {}^{T}F .
\]
Pour tordre $F$ par un $G$-torseur quand $G$ opère \emph{via} un morphisme $G
\to G'$, il suffit donc de tordre par le $G'$-torseur déduit $T'$ ; par
exemple, si $F$ porte quelques structures, on peut prendre $G' =
\mathrm{Aut}(F)$\footnote{La page définit ${}^{G'}E = E \times_G G'$ pour un
$G$-ensemble $E$ « à gauche », comme $G'$-ensemble « à g. », puis prend $F$
« à droite » ; les côtés ne s'accordent pas. Le produit $E \wedge_G G'$
demande $E$ à droite et donne un $G'$-ensemble à droite, ce qui est le cas
du torseur $T$ ; on a récrit l'article avec ces conventions. Le $G' =
\mathrm{Aut}(F)^{\circ}$ de la page 20 appelle la même remarque qu'à la
page 18. La fin de l'exemple, de « et $F^{T}$ » à « $G'$-torseur $T'$ », est
en partie illisible.}.

\textbf{(6) Produit des torseurs sous un groupe commutatif.} Si $G$ est
commutatif, la multiplication $G \times G \to G$ est un morphisme, et l'on
définit le produit de deux torseurs comme le torseur déduit par extension du
groupe :
\[
  X \wedge Y = (X \times Y) \wedge_{G \times G} G,
\]
et de même $(X_1 \times \cdots \times X_n) \wedge_{G^n} G$. \textbf{Un produit
de torseurs est un torseur}. Ce produit est associatif et commutatif à
isomorphisme canonique près, d'unité le torseur trivial $G_d$, et tout torseur
$X$ a un inverse : $X^{-1} = X \wedge_G G$, $G$ opérant sur lui-même par
$g \mapsto g^{-1}$ (un morphisme, $G$ étant commutatif), avec $X \wedge X^{-1}
\simeq G_d$ ; un tel inverse est unique à isomorphisme unique près compatible
aux trivialisations\footnote{La page note l'inverse additivement, $g \mapsto
-g$, et laisse la condition d'unicité en points de suspension. Les torseurs
sous un groupe commutatif forment ainsi ce qu'on appelle depuis Deligne
(SGA 4, exposé XVIII) une catégorie de Picard strictement commutative ; le
nom n'est pas sur la page. C'est cette structure qui sert, à partir de la
page 22, pour les torseurs sous $\{\pm 1\}$.}.

\textbf{(7)} Un article sur les orientations, qui commençait par
« $X \wedge Y$ », est biffé, puis un exemple encadré --- un ensemble fini $B$
muni d'une relation d'équivalence, $I = B/R$, et le produit
$\bigwedge_{i \in I} B_i$ --- est barré à son tour. Il annonce la page
suivante.

\pagerange{21}{21}
\section*{8. Structures cubiques (page 21)}

Soit $A$ un anneau dans lequel $2 \neq 0$, et $n \geq 1$. Soit $\mathcal{C}$
la catégorie des $A$-modules libres $E$ munis d'une décomposition en somme
directe $E = \bigoplus_{i \in I} E_i$ de $n$ sous-modules libres de rang $1$,
chaque $E_i$ muni d'une \emph{bi-base} $\{e_i, -e_i\}$ ; de façon
équivalente, des $A$-modules $E$ munis d'une partie $B$ à $2n$ éléments,
stable par $x \mapsto -x$, telle qu'un --- donc tout --- système de
représentants de $B/\{\pm 1\}$ soit une base. Soit $\mathcal{C}'$ la catégorie
des ensembles $B$ à $2n$ éléments munis d'une involution $\sigma$ sans point
fixe. Prenons pour morphismes, dans $\mathcal{C}$, les applications
$A$-linéaires $f$ avec $f(B) \subset B'$, et dans $\mathcal{C}'$ les
applications qui commutent à $\sigma$.

\textbf{Proposition.} Le foncteur évident $\mathcal{C} \to \mathcal{C}'$,
$(E, B) \mapsto (B, x \mapsto -x)$, est une équivalence de catégories. Un
quasi-inverse associe à $(B, \sigma)$, avec $I = B/\sigma$ et $B_i$ la classe
de $i$, le module
\[
  E = \bigoplus_{i \in I} E_i, \qquad E_i = B_i \wedge_{\{\pm 1\}} A,
\]
où $B_i$ est vu comme torseur sous $\{\pm 1\}$ et $\{\pm 1\}$ opère sur $A$
par homothéties ; $E_i$ est libre de rang $1$, de bi-base $\{[b, 1], [\sigma
b, 1] = -[b, 1]\}$\footnote{La formule de la page s'interrompt après le
signe $\wedge^{\{\pm 1\}}$ sur un symbole surchargé et barré, puis nomme
« $D$ », lecture, « le $A$-module $A_g$ [lecture] sur lequel $\pm 1$ opère
par homothéties » ; on écrit $A$. Le titre « Digression » est une lecture. La
page ne dit pas quels sont les morphismes ; avec ceux qu'on a choisis, le
foncteur est même bijectif sur les ensembles de morphismes, une application
linéaire étant déterminée par l'image d'une base. L'hypothèse $2 \neq 0$ est
exactement ce qu'il faut pour que $e_i \neq -e_i$, c'est-à-dire pour que $B$
ait $2n$ éléments.}.

Le groupe des automorphismes d'un objet est le groupe des permutations
signées $\{\pm 1\}^n \rtimes \mathfrak{S}_n$, le groupe des symétries du cube
de dimension $n$ --- d'où le titre. C'est le groupe hyperoctaédral, groupe de
Weyl des systèmes de racines de type $B_n$ et $C_n$ --- ceux-là mêmes que,
dans la lecture du dossier 88 (pages 2 à 4), le groupe de Weyl épinglé ne
distingue pas\footnote{Le nom et le rapprochement sont de nous.}. Le passage, numéroté
(8), prolonge l'exemple biffé de l'article (7) : $B$ muni de son involution est
un ensemble au-dessus de $I$ dont les fibres sont des torseurs sous $\{\pm
1\}$.

\pagerange{22}{23}
\section*{9. Un contour comme torseur tordu (pages 22 et 23)}

Soit $C$ un contour combinatoire connexe de longueur $n \geq 3$\footnote{La
page écrit « contour (combin.) convexe » ; « convexe », remplacé sans être
biffé par « (combin.) », n'a pas de sens pour un contour combinatoire, et la
suite demande la connexité.}. Il définit :
\begin{enumerate}
\item[1.] l'ensemble $\omega = \omega_C$ de ses deux orientations, torseur
sous $\{\pm 1\}$. On peut déjà en déduire le groupe des « translations » du
contour :
\[
  G^{+} = \mathrm{Aut}^{+}(C) \simeq \omega \wedge_{\{\pm 1\}} \mathbf{Z}_n,
\]
où $\{\pm 1\}$ opère sur $\mathbf{Z}_n$ par $x \mapsto -x$ ; l'isomorphisme
est $[o, k] \mapsto u_o^{k}$, bien défini parce que $u_{-o} = u_o^{-1}$
(page 7) ;
\item[2.] un torseur sous $G^{+}$ : l'ensemble de ses sommets.
\end{enumerate}
Inversement, un ensemble $\omega$ à deux éléments et un torseur $T$ sous
$\omega \wedge_{\{\pm 1\}} \mathbf{Z}_n$ définissent un contour : sommets $T$,
arêtes les paires $\{t, t + [o, 1]\}$, $o \in \omega$ ; et cela donne une
équivalence entre contours connexes de longueur $n$ et couples
$(\omega, T)$\footnote{La page énonce la réciproque sous le titre
« Exercice », biffé, et ne dit pas « équivalence » ; « (sinon compris) » est
une lecture. C'est, pour $\mathbf{Z}_n$, le groupe $\Gamma_\omega =
\mathbf{Z}/3\mathbf{Z} \wedge \omega$ que le dossier 89 (pages 2 à 6) écrit
pour $n = 3$.}. Le groupe d'automorphismes d'un couple $(\omega, T)$ est
$\mathbf{Z}_n \rtimes \{\pm 1\} = \mathbf{D}_n$, celui du contour : c'est la
forme intrinsèque de la page 18, où un contour était un torseur sous
$\mathbf{D}_n$.

\textbf{Exercice (DEA).} Soit $G = Z \rtimes \Gamma$, $\Gamma$ opérant sur
$Z$, avec $Z$ \emph{commutatif}. Se donner un $G$-torseur $T$ équivaut à se
donner
\begin{enumerate}
\item[a)] un $\Gamma$-torseur $\omega$ ;
\item[b)] un torseur $T_0$ sous le groupe tordu ${}^{\omega}Z = \omega
\wedge_\Gamma Z$.
\end{enumerate}
On passe de $T$ à $\omega = T \wedge_G \Gamma = T/Z$ et à $T_0 = T \wedge_G
(G/\Gamma) = T/\Gamma$, où $G/\Gamma$ est $Z$ muni de son action affine ;
c'est l'exercice dont l'énoncé sur les contours est le cas $Z = \mathbf{Z}_n$,
$\Gamma = \{\pm 1\}$, avec $T = \mathrm{Isom}(P_n, C)$, $\omega = \omega_C$,
$T_0 = S(C)$\footnote{L'hypothèse de commutativité n'est pas sur la page.
Sans elle, $T/\Gamma$ est encore un torseur, mais sous la forme tordue $T
\wedge_G Z$, $G$ opérant sur $Z$ par conjugaison, qui ne s'identifie à
$\omega \wedge_\Gamma Z$ qu'à automorphisme intérieur près --- le phénomène
de la page 2. Les deux lettres « $Z$ et $\Gamma$ », surchargées l'une sur
l'autre, ont un ordre incertain ; le sens n'en dépend pas, $Z$ étant le
facteur sur lequel $\Gamma$ opère.}.

\pagerange{24}{27}
\section*{10. Orientations et formules d'associativité (pages 24 à 27)}

\pagerange{24}{25}
\subsection*{Trois définitions d'une orientation (pages 24 et 25)}

Chaque fois, l'objet est un torseur $P$ sous un groupe $G$, muni d'un
caractère $s : G \to \{\pm 1\}$, et l'ensemble de ses orientations est le
$\{\pm 1\}$-torseur déduit
\[
  \omega = P \wedge_G \{\pm 1\} \quad (G \text{ opérant par } s),
\]
qui s'identifie à $P/\mathrm{Ker}(s)$ quand $s$ est surjectif\footnote{La page
écrit chaque fois « $\simeq P/\mathrm{Gl}^{+}$ », « $\simeq P/\mathrm{SO}$ »,
« $\simeq P/\mathfrak{A}_n$ » sans condition. Cela demande $s$ surjectif :
$n \geq 1$ en a) et b), $n \geq 2$ en c). Pour l'espace nul ou l'ensemble à au
plus un élément, $P$ est un point et $P/\mathrm{Ker}(s)$ aussi, tandis que
$P \wedge_G \{\pm 1\}$ a deux éléments ; c'est cette dernière définition qu'on
retient, et que demandent les formules d'associativité, où figurent des
ensembles vides et des singletons.}.
\begin{enumerate}
\item[a)] \emph{Espace vectoriel réel $E$ de dimension $n$} : $P =
\mathrm{Isom}(\mathbf{R}^n, E)$, l'ensemble des bases, sous $G = \mathrm{GL}(n,
\mathbf{R})$, et $s = \operatorname{sgn} \det$, de noyau $\mathrm{GL}^{+}(n,
\mathbf{R})$.
\item[b)] \emph{Module quadratique libre $E$ de rang $n$ sur un anneau
commutatif $k$ où $2$ est inversible, possédant une base orthonormale} : $P$
l'ensemble des bases orthonormales, sous $G = \mathrm{O}(n, k)$, et $s =
\det$, de noyau $\mathrm{SO}(n, k)$. Il faut supposer $k$ connexe --- sans
idempotent autre que $0$ et $1$, par exemple local ou intègre ---, sans quoi
$\det$ prend ses valeurs dans le groupe $\mu_2(k)$ des racines carrées de
l'unité, plus gros que $\{\pm 1\}$\footnote{Hypothèse ajoutée. Si $x^2 = 1$
et $2$ est inversible, $(1 + x)/2$ est idempotent ; c'est donc exactement la
connexité qui assure $\mu_2(k) = \{\pm 1\}$. Sans elle, l'ensemble des
orientations est un torseur sous $\mu_2(k)$. « Quadratiques » est ajouté
au-dessus de la ligne.}.
\item[c)] \emph{Ensemble fini $I$ de cardinal $n$} : $P =
\mathrm{Isom}(\{1, \ldots, n\}, I)$, l'ensemble des numérotations, sous
$\mathfrak{S}_n$, et $s$ la signature, de noyau $\mathfrak{A}_n$.
\end{enumerate}

\textbf{Compatibilités.} Elles résultent toutes de la transitivité de
l'extension du groupe structural (article (5)) : si $G_1 \to G_2$ est un
morphisme compatible aux caractères, et $P_2 = P_1 \wedge_{G_1} G_2$, les
orientations définies par $P_1$ et $P_2$ coïncident.
\begin{enumerate}
\item[a)--b)] Pour $k = \mathbf{R}$ : $\mathrm{O}(n, \mathbf{R}) \subset
\mathrm{GL}(n, \mathbf{R})$, et $\operatorname{sgn}\det$ y induit $\det$ ; les
bases orthonormales engendrent toutes les bases, d'où $\omega_E$ calculé dans
a) $=$ $\omega_E$ calculé dans b)\footnote{La page écrit le diagramme pour
$k$ quelconque, avec $\operatorname{sgn} \det$ sur $\mathrm{Gl}(n, k)$, qui n'a
de sens que pour $k = \mathbf{R}$, ou un corps ordonné.}.
\item[b)--c)] $k^I$, muni du produit scalaire $\sum_{i} x_i y_i$, est un module
quadratique $E(I, k)$ dont la base canonique est orthonormale ; une
numérotation $i_1, \ldots, i_n$ de $I$ donne la base orthonormale ordonnée
$e_{i_1}, \ldots, e_{i_n}$, et $\mathfrak{S}_n \to \mathrm{O}(n, k)$ (matrices
de permutation) vérifie $\det = \operatorname{sgn}$. D'où $\omega_{E(I, k)}
\simeq \omega_I$, et en particulier $\omega_{E(I, \mathbf{R})} \simeq
\omega_I$\footnote{« Ord. » dans la marge est une lecture. La dernière ligne
de la page 25, sous la formule, est très incertaine (« \ldots{} considéré
\ldots{} vect. réel se pense ») et n'est pas reprise.}.
\end{enumerate}

\pagerange{26}{27}
\subsection*{Formules d'associativité (pages 26 et 27)}

\textbf{a) Espaces vectoriels réels.} Soit $(E_i)_{i \in I}$ une famille finie
d'espaces vectoriels réels de dimension finie $d_i$, $E = \bigoplus_i E_i$, et
$I^{+}$ (resp.\ $I^{-}$) l'ensemble des $i$ tels que $d_i$ soit pair (resp.\
impair). On a un isomorphisme canonique de $\{\pm 1\}$-torseurs
\[
  \boxed{\ \omega_E \simeq \omega_{I^{-}} \wedge \bigwedge_{i \in I} \omega_{E_i}\ }
\]
\emph{Construction.} Une numérotation $(i_1, \ldots, i_{n^-})$ de $I^{-}$, une
numérotation $(j_1, \ldots, j_{n^+})$ de $I^{+}$ et, pour tout $i$, une base
ordonnée $B_i = (e_{i,1}, \ldots, e_{i,d_i})$ de $E_i$ donnent, par
concaténation des $B_i$ dans l'ordre $i_1, \ldots, i_{n^-}, j_1, \ldots,
j_{n^+}$, une base de $E$, donc un élément de $\omega_E$. On vérifie :
\begin{itemize}
\item changer un $B_i$ par une matrice de déterminant positif change la base
de $E$ par une matrice de déterminant positif, et par un déterminant négatif,
par un déterminant négatif : l'élément de $\omega_E$ ne dépend que des
orientations des $E_i$, et change de signe avec chacune ;
\item permuter $I^{-}$ par $\sigma^{-}$ change la base par une matrice de
déterminant $\operatorname{sgn}\sigma^{-}$ --- échanger deux blocs consécutifs
de tailles $d$ et $d'$ coûte $(-1)^{d d'}$, soit $-1$ pour deux blocs
impairs ;
\item permuter $I^{+}$ ne change rien, chaque échange mettant en jeu un bloc
pair.
\end{itemize}
L'application $\mathrm{numérot}(I^-) \times \mathrm{numérot}(I^+) \times
\prod_i \mathrm{bases}(E_i) \to \omega_E$ passe donc au quotient en une
application $\omega_{I^-} \times \prod_i \omega_{E_i} \to \omega_E$ compatible
à l'action de $\{\pm 1\}$ sur chaque facteur, c'est-à-dire en l'isomorphisme
encadré. Le même argument montre que l'on peut concaténer selon n'importe quel
ordre total sur $I$ induisant l'ordre donné sur $I^{-}$, comme le note la
marge\footnote{La page écrit la base concaténée sur les seuls indices $i_1,
\ldots, i_n$ de $I^-$, alors que la numérotation de $I^+$, ajoutée en
interligne, doit y entrer aussi. « Quotient » et « donc définit le même él. »
sont des lectures. Le calcul par $(-1)^{dd'}$ est le nôtre ; la page traite
une transposition de deux voisins $\alpha$, $\alpha + 1$ et s'arrête sur
« la matrice de passage est de dét. $+1$ en tous cas… ». C'est la règle de
signe de Koszul, celle qui gouverne la droite déterminant graduée de Knudsen
et Mumford (1976) ; les noms ne sont pas sur la page.}.

\textbf{b) Modules quadratiques.} Même formule, même démonstration, pour une
famille finie de modules quadratiques libres à base orthonormale, sur $k$
comme en b) de la page 24.

\textbf{c) Ensembles finis.} Soit $p : J \to I$ une application d'ensembles
finis --- par exemple la projection de $J$ sur son quotient par une relation
d'équivalence ---, $J_i = p^{-1}(i)$ et $I^{-}$ l'ensemble des $i$ dont la
fibre a un cardinal impair. Alors
\[
  \omega_J \simeq \omega_{I^{-}} \wedge \bigwedge_{i \in I} \omega_{J_i},
\]
ce qui est a) appliqué à $\mathbf{R}^J = \bigoplus_i \mathbf{R}^{J_i}$, compte
tenu de $\omega_{E(J, \mathbf{R})} \simeq \omega_J$\footnote{« Surjection si
on veut » est une lecture, et l'indice $J_i$ est empâté par une tache
d'encre. Une fibre vide est paire et a pour orientations le torseur trivial,
d'où l'indifférence. La déduction à partir de a) est de nous ; c'est l'une
des compatibilités que la page se donne à vérifier.}. La page se donne
ensuite à établir la « compatibilité entre les isomorphismes canoniques a),
b), c) », sans l'écrire.

\textbf{Cas des extensions.} Pour un drapeau $0 = E_0 \subset E_1 \subset
\cdots \subset E_d = E$, la page demande de « définir un isomorphisme
canonique »
\[
  \omega_E \simeq \bigwedge_{1 \leq i \leq d} \omega_{E_i/E_{i-1}} .
\]
Il s'obtient en relevant des bases des quotients et en les concaténant dans
l'ordre du drapeau ; le résultat ne dépend pas des relèvements, qui changent
la base par une matrice unipotente triangulaire par blocs. Il n'y a pas de
terme correctif : l'ordre est donné par le drapeau\footnote{La démonstration
est de nous ; la page s'arrête sur l'énoncé.}.

\pagerange{28}{32}
\section*{11. La lettre et les deux annonces de cours (pages 28 à 32)}

Les pages 28 et 29 sont deux copies --- un carbone pâle et une photocopie ---
d'une lettre dactylographiée qu'il signe à Villecun le 30 janvier 1978, adressée
au ministre de l'Éducation au sujet de l'ordonnance du 2 novembre 1945. C'est la
« lettre (1978) » de l'inventaire ; elle ne porte pas de mathématiques, la
transcription ne la reproduit pas, et on n'en dit pas plus ici.

\pagerange{30}{31}
\subsection*{Cours C4 1977--1978 : géométrie de l'icosaèdre (pages 30 et 31)}

L'annonce, dactylographiée sur papier à en-tête de l'Institut de
mathématiques, part d'un étonnement : l'icosaèdre, étudié depuis plus de deux
mille ans, ne semble jamais avoir été étudié systématiquement du point de vue
combinatoire, ni de celui de ses réalisations géométriques sur des corps,
voire des anneaux, quelconques. Le cours se veut une « investigation commune »
partant de l'intuition combinatoire la plus naïve, sur des modèles en carton
dont les arêtes sont coloriées.

Le programme : expliciter le dictionnaire entre la géométrie combinatoire de
l'icosaèdre orienté et celle de l'ensemble « orienté » à cinq éléments --- les
cinq couleurs des arêtes ---, et étudier le groupe d'automorphismes commun aux
deux situations, le groupe alterné $\mathfrak{A}_5$, plus petit groupe simple
non abélien\footnote{Les trente arêtes de l'icosaèdre se répartissent en
quinze paires d'arêtes opposées, et celles-ci en cinq triples de paires deux à
deux orthogonales ; le groupe des rotations permute ces cinq triples, et c'est
ainsi qu'il s'identifie à $\mathfrak{A}_5$. L'annonce ne dit pas quel
coloriage elle a en vue ; celui-ci est le classique.} ; puis construire des
réalisations de l'icosaèdre combinatoire comme polyèdre régulier sur les réels
et sur d'autres corps, ce qui obligera à faire de la géométrie vectorielle,
affine et projective en caractéristique quelconque. Les caractéristiques deux
et cinq y jouent un rôle particulier, et l'étude combinatoire s'y exprime par
la géométrie sur $\mathbf{F}_4$ et $\mathbf{F}_5$, à travers les isomorphismes
exceptionnels
\begin{align*}
  \mathfrak{A}_5 &\simeq \mathrm{SL}(2, \mathbf{F}_4) \simeq \mathrm{PGL}(2,
    \mathbf{F}_4) \simeq \mathrm{SO}(3, \mathbf{F}_4), \\
  \mathfrak{S}_5 &\simeq \mathrm{PGL}(2, \mathbf{F}_5) \simeq \mathrm{SO}(3,
    \mathbf{F}_5),
\end{align*}
dont la page 13 contenait déjà les deux premiers\footnote{La frappe écrit
$\mathrm{Sl}$ et $\mathrm{GP}(1, \cdot)$ ; les lettres gothiques y sont
ajoutées à la main. En caractéristique $2$, une forme quadratique non
dégénérée en dimension $3$ a une forme bilinéaire associée de rang $2$, et le
groupe orthogonal, où le déterminant vaut toujours $1$, est isomorphe à
$\mathrm{Sp}(2, \mathbf{F}_4) = \mathrm{SL}(2, \mathbf{F}_4)$ ;
$\mathrm{SO}(3, \mathbf{F}_4)$ doit s'entendre ainsi. Pour $q$ impair,
$\mathrm{SO}(3, \mathbf{F}_q) \simeq \mathrm{PGL}(2, \mathbf{F}_q)$.}.

\pagerange{32}{32}
\subsection*{Cours DEA 1977--1978 : géométrie arithmétique des polygones réguliers (page 32)}

Le second cours, conçu lui aussi comme une investigation en commun, prend
pour objet la notion de polygone régulier au-dessus d'un corps de base, puis
d'un anneau de base quelconque. Il doit familiariser avec la géométrie
vectorielle, affine et projective, avec ou sans métrique, en dimensions un,
deux et trois, et avec les formes imprévues que prennent en caractéristique
positive, notamment sur les corps finis, des notions familières comme celle de
réflexion ou de centre d'un polygone régulier. Le besoin d'une géométrie sur
un anneau de base général --- tel l'anneau « sur lequel s'écrit » le $n$-gone
régulier « le plus général » --- doit mener à quelques points de fondements de
géométrie algébrique ; si le temps le permet, on tentera de construire
l'icosaèdre régulier « général » sur un anneau de base fini sur $\mathbf{Z}$,
ayant des corps résiduels de toutes caractéristiques, la caractéristique deux
promettant « du fil à retordre ». L'annonce insiste sur la participation
active et le travail personnel de l'étudiant, et conseille de suivre en même
temps le cours sur l'icosaèdre\footnote{Au crayon, un « ! » dans la marge en
face de la première ligne et un trait sous « en commun » ; rien n'indique de
quelle main.}.

Les feuillets manuscrits des pages 2 à 27 sont l'outillage de ces deux cours :
le polygone combinatoire comme torseur sous $\mathbf{D}_n$, son enveloppe
canonique, ses orientations, les petits groupes comme groupes de géométrie
finie. Les mêmes objets se retrouvent ailleurs dans le groupe : les triangles,
hexagones et polygones combinatoires du dossier 89, les polygones
« bi-circulés » du dossier 81 (pages 40 à 52), et, dans le langage des
cartes et du groupe cartographique, le dossier 88.

\end{document}
